% Additive research section.  No source manuscript has been modified.
% Requires amsmath, amssymb, amsthm; theorem, proposition, corollary,
% remark environments.  Citation keys are left to the integrating article.
\section{All spectral coefficients from ordinary logarithmic Mellin moments}
\label{mellin:section}

The incoming \emph{Differential Polygamma and Lerch Jets} report
\cite[Section~9]{incoming-differential} asks for
spatial functionals selecting higher coefficients of its residue Dirichlet
germ, with the exact logarithmic normalization and endpoint constants.
The construction below answers that question at every order.  The
Mittag--Leffler expansion and the scalar residue filter are inherited
inputs; the shift-difference Mellin transform, its ordinary subtractions,
and the higher-coefficient formulas are the continuation proved here.
This is a claim of progress relative to the inspected material, not a
historical priority claim for all specializations of Mellin calculus.

\subsection{The data and the spatial normalization}

Let $P$ be a polynomial in finitely many variables $u_0,u_1,\ldots$ and
write
\[
 f_P(x)=P(\psi(x),\psi'(x),\ldots),\qquad
 \phi_P(Y)=P(Y,0,0,\ldots).
\]
Let its principal part at $-n$ be
\[
 \sum_{j=1}^{M}\frac{p_j(n)}{(x+n)^j},\qquad
 \Delta p_j(n)=p_j(n)-p_j(n-1),\qquad
 P_0(x)=\sum_{j=1}^{M}p_j(0)x^{-j}.
\]
Zero coefficients may be included.  These finite Laurent coefficients
are the harmonic polynomials supplied by the incoming report.  In
particular, $p_j(n)=O((1+\log n)^d)$ and
$\Delta p_j(n)=O((1+\log n)^d/n)$ for some fixed $d$.
Define
\begin{align}
 F_P(x)&=f_P(x)-f_P(x+1)-P_0(x),\label{mellin:F}\\
 c_P(L)&=\phi_P'(L)+p_1(0)=\sum_{r=0}^{R}c_rL^r,\label{mellin:c}\\
 h_P(x)&=F_P(x)+\mathbf 1_{[1,\infty)}(x)\frac{c_P(\log x)}x.
 \label{mellin:h}
\end{align}
The cutoff is exactly $1$, the same scale as the finite-part integral
on $[0,1]$.  A value assigned to the step function at $1$ has no effect.

\begin{theorem}[Ordinary Mellin realization of every residue coefficient]
\label{mellin:main}
The function $F_P$ extends analytically through $x=0$, and
\[
 h_P(x)=O((1+\log x)^d/x^2)\quad(x\longrightarrow\infty).
\]
Consequently
\begin{equation}\label{mellin:A}
 A_P(u)=\int_0^\infty x^{-u}h_P(x)\,dx
\end{equation}
is holomorphic in $-1<\operatorname{Re}u<1$.  Its derivatives there
are ordinary absolutely convergent integrals.  Put
\begin{equation}\label{mellin:B}
 \mathcal B_P(u)=\frac{\sin\pi u}{\pi}
 \left(A_P(u)-\sum_{r=0}^{R}\frac{c_r r!}{u^{r+1}}\right).
\end{equation}
This is a meromorphic germ at zero.  If
\[
 D_j(v)=\sum_{n\ge1}\frac{\Delta p_j(n)}{n^v},\qquad D_P=D_1,
\]
then the exact meromorphic identity is
\begin{equation}\label{mellin:bridge}
 \boxed{\mathcal B_P(u)=D_P(u)+
 \sum_{j=2}^{M}\frac{(u)_{j-1}}{(j-1)!}D_j(u+j-1).}
\end{equation}
Every series in the sum over $j\ge2$ is normally convergent on
$\operatorname{Re}u>-1$.  In the original strip
$0<\operatorname{Re}u<1$ one may equivalently write
\begin{equation}\label{mellin:unsubtracted}
 \mathcal B_P(u)=\frac{\sin\pi u}{\pi}
                  \int_0^\infty x^{-u}F_P(x)\,dx.
\end{equation}
\end{theorem}

\begin{proof}
Subtract the normalized Mittag--Leffler formula for $f_P(x+1)$ from
that for $f_P(x)$.  The constants cancel.  Reindexing a finite cutoff
and passing to its limit gives
\begin{equation}\label{mellin:partialfractions}
 F_P(x)=\sum_{n\ge1}\sum_{j=1}^{M}
               \frac{\Delta p_j(n)}{(x+n)^j}.
\end{equation}
Indeed the uncancelled last term is a principal part evaluated at
$x+N+1$, and tends to zero as
$O((1+\log N)^d/N)$.  The new series is normally convergent on
compact sets away from the negative integers.  In particular it is
analytic at zero.  This also follows directly from
$f_P-P_0$ and $f_P(x+1)$ being analytic there.

On the positive real axis, the usual differentiated digamma asymptotics
give
\[
 f_P'(x)=\frac{\phi_P'(\log x)}x
           +O((1+\log x)^d/x^2).
\]
Every monomial containing a positive polygamma derivative contributes
only to the error.  Integrating this estimate from $x$ to $x+1$, and
using $P_0(x)=p_1(0)/x+O(x^{-2})$, proves
\[
 F_P(x)=-\frac{c_P(\log x)}x
              +O((1+\log x)^d/x^2).
\]
Thus the subtraction in \eqref{mellin:h} proves all assertions about
$A_P$, including differentiation under the integral sign.

For $0<\operatorname{Re}u<1$ absolute convergence permits insertion
of \eqref{mellin:partialfractions}.  The elementary beta integral is
\[
 \int_0^\infty\frac{x^{-u}}{(x+n)^j}\,dx
 =n^{1-u-j}\frac{\Gamma(1-u)\Gamma(j-1+u)}{\Gamma(j)}.
\]
Its multiplication by $\sin(\pi u)/\pi$ is
$n^{1-u-j}(u)_{j-1}/(j-1)!$, including $j=1$.
The coefficient bound above makes the resulting sum absolutely
convergent.  Also
\[
 \int_1^\infty x^{-u-1}(\log x)^r\,dx=\frac{r!}{u^{r+1}}.
\]
These two computations prove \eqref{mellin:bridge} and
\eqref{mellin:unsubtracted} in the strip.  The $j\ge2$ series
converge locally normally for $\operatorname{Re}u>-1$; the scalar
residue filter continues $D_P$ meromorphically there.  Uniqueness of
continuation proves the germ identity.  No integral at a divergent
endpoint has been differentiated without its stated subtraction.
\end{proof}

\subsection{An explicit coefficient formula at every order}

For $m,r\ge0$ define the finite polynomial
\begin{equation}\label{mellin:E}
 E_{m,r}(L)=[u^r]e^{-Lu}\prod_{k=1}^{m}(1+u/k).
\end{equation}
Here an empty product is one.  The first three are
\[
 E_{m,0}=1,\qquad E_{m,1}=H_m-L,\qquad
 E_{m,2}=\frac{(H_m-L)^2-H_m^{(2)}}2.
\]
Thus these are finite harmonic polynomials, or equivalently Bell
polynomials with cumulants $H_m-L$ and
$(-1)^{k-1}(k-1)!H_m^{(k)}$ for $k\ge2$.

Write $\mu_k(P)=\int_0^\infty h_P(x)(\log x)^k\,dx$.
For $q\ge0$ put
\begin{align}
 B_q(P)={}&
 \sum_{\substack{k\ge0\\2k+1\le q}}
 \frac{(-1)^{k+q-2k-1}\pi^{2k}}
 {(2k+1)!(q-2k-1)!}\,\mu_{q-2k-1}(P)\notag\\
 &-\sum_{\substack{0\le r\le R\\q+r\ {
m even}}}
 c_r r!\frac{(-1)^{(q+r)/2}\pi^{q+r}}{(q+r+1)!}.
 \label{mellin:Bq}
\end{align}
The first sum is empty for $q=0$.  It is often clearer computationally
to take the coefficient directly in \eqref{mellin:B}; formula
\eqref{mellin:Bq} specifies every sign and factorial.

\begin{corollary}[Higher coefficients and the complete Stieltjes term]
\label{mellin:coefficients}
For every integer $q\ge1$,
\begin{equation}\label{mellin:Dq}
 \boxed{[u^q]D_P(u)=B_q(P)-
 \sum_{j=2}^{M}\sum_{n\ge1}
 \frac{\Delta p_j(n)}{(j-1)n^{j-1}}
 E_{j-2,q-1}(\log n).}
\end{equation}
The double sum is finite in $j$ and absolutely convergent in $n$.
If the inherited scalar residue polynomial is
\[
 R_P(L)=-\frac{\phi_P(L+i\pi)-\phi_P(L-i\pi)}{2i\pi},
 \qquad R_P'(L)=\sum_{r=0}^{S}a_rL^r,
\]
put
$\delta_n=\Delta p_1(n)-R_P'(\log n)/n$.
For every $q\ge0$ its entire ordinary-series counterpart is
\begin{equation}\label{mellin:stieltjes}
 \boxed{[u^q]D_P(u)=\frac{(-1)^q}{q!}
 \left(\sum_{r=0}^{S}a_r\gamma_{q+r}
                  +\sum_{n\ge1}\delta_n(\log n)^q\right).}
\end{equation}
The second series converges absolutely.  In particular, the same
scalar projection controls \emph{all} Stieltjes corrections.  If
$\phi_P=0$, every displayed $\gamma_{q+r}$ term vanishes.
\end{corollary}

\begin{proof}
Expanding $A_P$ in \eqref{mellin:B} gives \eqref{mellin:Bq}.
For $j\ge2$,
\[
 \frac{(u)_{j-1}}{(j-1)!}n^{-u}
 =\frac{u}{j-1}e^{-u\log n}\prod_{k=1}^{j-2}(1+u/k).
\]
Taking the coefficient of $u^q$ in the normally convergent identity
\eqref{mellin:bridge} proves \eqref{mellin:Dq}.  The logarithmic
factor grows by only a fixed power, so the convergence assertion
persists at every $q$.

The scalar filter gives
$\delta_n=O((1+\log n)^d/n^2)$ and
\[
 D_P(u)=\sum_r a_r(-1)^r\zeta^{(r)}(1+u)
                         +\sum_{n\ge1}\delta_nn^{-u}.
\]
With the convention
$\zeta(1+u)=u^{-1}+\sum_{m\ge0}(-1)^m\gamma_m u^m/m!$,
the coefficient of $u^q$ in
$(-1)^r\zeta^{(r)}(1+u)$ is
$(-1)^q\gamma_{q+r}/q!$.  This proves \eqref{mellin:stieltjes}.
\end{proof}

For $q=1$, the last sum in \eqref{mellin:Dq} is exactly the inherited
higher-pole correction
\[
 J_P=\sum_{j=2}^{M}\sum_{n\ge1}
              \frac{\Delta p_j(n)}{(j-1)n^{j-1}}.
\]
Hence $[u]\mathcal B_P=J_P+[u]D_P$, and the prior finite-part theorem
is recovered with its original endpoint constant:
\[
 \operatorname{FP}\int_0^1f_P(x)\,dx
       =C_P+[u]\mathcal B_P(u),\qquad
 C_P=\frac1{2\pi}\int_{-\pi}^{\pi}\phi_P(i\theta)\,d\theta.
\]
This is a normalization check, not a new proof of that antecedent.

\begin{corollary}[The first new logarithmic spatial functional]
\label{mellin:mixedq2}
If $P$ is mixed, meaning $\phi_P=0$, then
\begin{equation}\label{mellin:q2}
 \boxed{[u^2]D_P(u)=
 -\int_0^\infty h_P(x)\log x\,dx+\frac{\pi^2p_1(0)}6
 -\sum_{j=2}^{M}\sum_{n\ge1}
 \frac{\Delta p_j(n)(H_{j-2}-\log n)}{(j-1)n^{j-1}}.}
\end{equation}
All integrals and series are ordinary and absolutely convergent.
\end{corollary}

The local coefficient $p_1(0)$ is essential even in the mixed case.
It is the correction from $\sin(\pi u)/(\pi u)$, and its coefficient
is invisible at first order.  Thus disappearance of the Stieltjes
terms does not permit omission of this second-order endpoint term.

\subsection{Concrete polygamma and nonlinear examples}

For $P=u_1^2$ one has
\[
 p_4(n)=1,\quad p_2(n)=2(H_n^{(2)}+\zeta(2)),\quad
 p_1(n)=4(H_n^{(3)}-\zeta(3)).
\]
All omitted principal-part coefficients vanish.  Therefore
\begin{equation}\label{mellin:trigammasquare}
 \mathcal B_{u_1^2}(u)=(4+2u)\zeta(u+3).
\end{equation}
Let $Q(x)=\psi'(1+x)$.  An ordinary all-order generating formula is
obtained from \eqref{mellin:B} using
\[
 h_{u_1^2}(x)=
 \begin{cases}
 2\{Q(x)-\zeta(2)+2\zeta(3)x\}/x^2,&0<x<1,\\
 2\{Q(x)-\zeta(2)\}/x^2,&x>1.
 \end{cases}
\]
Its quadratic coefficient yields the explicit identity
\begin{equation}\label{mellin:trigamma-log}
 \boxed{\begin{aligned}
 &\int_0^1\frac{\psi'(1+x)-\zeta(2)+2\zeta(3)x}{x^2}\log x\,dx\\
 &\quad+\int_1^\infty\frac{\psi'(1+x)-\zeta(2)}{x^2}\log x\,dx\\
 &\qquad=-\zeta''(3)-\zeta'(3)-\frac{\pi^2\zeta(3)}3.
 \end{aligned}}
\end{equation}
Primes on $\zeta$ denote derivatives in the spectral order.  The
linear polygamma example provides a particularly simple independent
check of the endpoint term in \eqref{mellin:q2}.

For an example retaining a nonlinear spatial integrand, take $P=u_1^3$.
Write $Z(s,t)=\sum_{n>m\ge1}n^{-s}m^{-t}$ in its convergent region.
Finite Laurent extraction gives
\begin{align*}
 p_6&=1,&p_5&=0,&p_4&=3c_2,&p_3&=6c_3,\\
 p_2&=9c_4+3c_2^2,&p_1&=12c_5+12c_2c_3,
 \qquad c_r=H_n^{(r)}+(-1)^r\zeta(r).
\end{align*}
Inserting their exact differences into \eqref{mellin:bridge} proves
the finite depth-two identity
\begin{equation}\label{mellin:trigammacube}
 \boxed{\begin{aligned}
 \mathcal B_{u_1^3}(u)={}&
 \frac{u^3+9u^2+32u+48}{2}\,\zeta(u+5)\\
 &+6(u+2)\{Z(u+3,2)+\zeta(2)\zeta(u+3)\}\\
 &+12\{Z(u+2,3)-\zeta(3)\zeta(u+2)\}.
 \end{aligned}}
\end{equation}
All $Z$ functions in a neighborhood of zero are ordinary absolutely
convergent strict series.  Arbitrary higher order derivatives may be
taken there; no identity restricted to integer arguments is being
differentiated.

For clarity, put $K=\zeta(5)+\zeta(2)\zeta(3)$ and
\begin{align*}
 U(x)={}&\frac{3Q(x)}{x^4}+\frac{3Q(x)^2}{x^2}
 -\frac{3\zeta(2)}{x^4}+\frac{6\zeta(3)}{x^3}
 -\frac{9\zeta(4)+3\zeta(2)^2}{x^2}.
\end{align*}
The ordinary function for \eqref{mellin:A} is
\[
 h_{u_1^3}(x)=U(x)+\frac{12K}{x}\mathbf1_{(0,1)}(x).
\]
The displayed cancellations make it bounded at zero; for numerical
work its convergent Taylor series avoids loss of digits there.
With $\partial_1$ denoting a genuine first-order spectral derivative,
the second coefficient of \eqref{mellin:trigammacube} is
\begin{align}
 B_2(u_1^3)={}&12\zeta''(5)+16\zeta'(5)+\tfrac92\zeta(5)
 +6(\partial_1^2+\partial_1)Z(3,2)\notag\\
 &+6\zeta(2)\{\zeta''(3)+\zeta'(3)\}
 +6\partial_1^2 Z(2,3)-6\zeta(3)\zeta''(2).
 \label{mellin:cube-B2}
\end{align}
It follows that the following nonlinear logarithmic integral has a
finite, convergent depth-two evaluation:
\begin{equation}\label{mellin:cube-log}
 \boxed{\int_0^\infty h_{u_1^3}(x)\log x\,dx
                  =-B_2(u_1^3)-2\pi^2K.}
\end{equation}
This evaluates the logarithmic spatial functional in a specified
spectral family.  It makes no claim that the retained derivatives
of $Z$ are arithmetically irreducible.

\subsection{What this settles and what remains}

The higher-spectral-coefficient question is resolved for every
differential polynomial: the spatial kernel, the scale, all endpoint
terms, the higher-pole corrections, and every Stieltjes coefficient are
explicit.  The inherited convergent-word compiler applies to the
coefficients in \eqref{mellin:Dq}: it only acquires powers of
$\log n$, hence higher derivatives in the existing spectral arguments.
Thus these formulas introduce no additional kind of multiple-zeta
coordinate.  They do not settle relations between the resulting
evaluated coordinates.  A separate next problem is to extend the
construction to independent shifts, where the pole lattices no longer
align under the unit difference, and to characterize which linear
combinations eliminate all higher-depth spectral derivatives.
